\documentclass[12pt]{article}
\usepackage[utf8]{inputenc}
\usepackage[T1]{fontenc}
\usepackage{hyperref}
\title{Kwamera Language Resources for K-12 Curriculum Development}
\author{Compiled by Claude Sonnet 5 (claude-sonnet-5) for Andr\'as Kornai}
\date{}
\begin{document}
\maketitle

\section{Introduction}
Kwamera (South Tanna; endonym Nafe) is an Oceanic language spoken on the southeastern coast of Tanna Island, Vanuatu, by about 3,500 people. It is taught as a subject in some schools. The searches found one printed dictionary, one online talking dictionary, and several general reference pages. I found no standalone Kwamera grammar.

\section{Primary Educational Resources}
No school primers or textbooks turned up in the searches. The orthography was revised by a Tanna orthographic committee in the 1980s, after missionaries first devised it. This is relevant to any materials produced (see Omniglot).
\begin{itemize}
\item Omniglot, Kwamera: \url{https://www.omniglot.com/writing/kwamera.htm}
\end{itemize}

\section{Linguistic and Reference Materials}
\begin{itemize}
\item Lamont Lindstrom, \emph{Kwamera Dictionary} (Pacific Linguistics C-95, 1986), library record: \url{https://find.library.upenn.edu/catalog/995510553503681}
\item Wikipedia, Kwamera language: \url{https://en.wikipedia.org/wiki/Kwamera_language}
\item Languages of Tanna (N. Thieberger): \url{https://nthieberger.net/vanlangs/TANNA.HTM}
\item Cambridge Language Centre, Tanna languages: \url{https://www.langcen.cam.ac.uk/resources/langt/tanna.html}
\end{itemize}

\section{Religious and Community Resources}
None found in the searches.

\section{Digital and Audio Resources}
\begin{itemize}
\item Capell's archival notes on Tanna languages, including Kwamera (PARADISEC): \url{https://www.paradisec.org.au/fieldnotes/VTAN.htm}
\item Verbix Kwamera verb page (low authority): \url{https://docs.verbix.com/Languages/Kwamera}
\end{itemize}
The Nafe (Kwamera) Talking Dictionary from Swarthmore College (K. D. Harrison and M. Wahe) has about 3,233 entries with audio, images and video. It was mentioned in search results, but I did not see its URL, so it is not listed with a link.

\section{Curriculum Development Context}
The language is still used by adults, and it is taught in some schools. The Lindstrom dictionary and the talking dictionary are the strongest foundations for a curriculum. A grammar for teachers would be a gap to fill. No named local contacts were found.

\end{document}